% Stat 220 syllabus, Fall 2026.
% Source file. Edit this, then compile with:
%     pdflatex 220_Syllabus_F2026.tex
% Structure adapted from the F2025 syllabus (F2026/Syllabus/220_Syllabus.tex),
% which is kept unchanged for reference.

\documentclass[11pt]{article}
\usepackage[margin=1in]{geometry}
\usepackage[pdftex]{graphicx}
\usepackage{multirow}
\usepackage{setspace}
\usepackage{enumitem}
\pagestyle{plain}
\setlength\parindent{0pt}
\usepackage{hyperref}

\begin{document}

% ---------------------------------------------------------------- course info
\begin{tabular}{ l l }
  \multirow{3}{*}{} & \LARGE Stat 220 \\
  & \LARGE Statistical Modeling for Data Science \\
  & \LARGE Fall 2026, MWF 9:00--9:50, 1102 JKB \\
\end{tabular}
\vspace{8mm}

\begin{tabular}{ l l }
  \multirow{6}{*}{} & \large Professor: Robert Richardson \\
  & \hspace{5mm} \large email: richardson@stat.byu.edu \\
  & \hspace{5mm} \large Office: 2186 WVB \\
  & \hspace{5mm} \large Open Zoom Room: 8014223736 \\
  & \hspace{5mm} \large Office Hours: Mondays 1:15--2:15 and Thursdays 9:00--10:00, 2186 WVB \\
  & \hspace{5mm} \large 801-422-3736 \\
\end{tabular}
\vspace{5mm}

\begin{tabular}{ l l }
  \multirow{6}{*}{} & \large TAs: \\
  & \hspace{5mm} \large Phoenix Nettesheim, phnetty@byu.edu \\
  & \hspace{5mm} \large Joseph Asay, jasay1@byu.edu \\
  & \hspace{5mm} \large Averi Carlson, averipc@byu.edu \\
  & \hspace{5mm} \large Gordon Lam, klam20@byu.edu \\
  & \hspace{5mm} \large TA office hours will be posted on Learning Suite. \\
\end{tabular}
\vspace{5mm}

% ------------------------------------------------------------- course details
\section*{Course Details}

\textbf{\large Course Description:}

This course is about what you can actually conclude from data, and what you cannot.

Computation is cheap and getting cheaper. Anyone can fit a model in three lines. The skill worth
having is the judgment that turns a number into a claim you can defend: knowing what a method
assumes, what happens when those assumptions break, and how to tell a real signal from a
coincidence that looks like one.

So each unit does two things. It builds a piece of the statistical toolkit, and then it spends
just as much time on how that piece fails. You will fit models in this class. The harder question,
and the one we keep returning to, is whether the model has earned the conclusion you want to draw
from it.

Most students here are studying computer science rather than statistics or mathematics. The course
is demanding, but very little is assumed. Concepts get built up rather than asserted, and the code
companions let you watch a mechanism move instead of taking it on faith.

~

\textbf{Prerequisite(s):} Stat 121 or Stat 201; CS 110 or CS 111

~

\textbf{\large Texts:} No textbook is required.

\begin{itemize}[topsep=2pt,itemsep=1pt]
\item (Optional) \textbf{Foundations of Statistics for Data Scientists: With R and Python} by
  Alan Agresti and Maria Kateri. This is the closest published book to the course. I will point
  you to sections of it for further study. Its examples are in R, but there is a Python supplement
  at \url{https://stat4ds.rwth-aachen.de/pdf/DS_Python_webAppendix.pdf} and data sets at
  \url{https://stat4ds.rwth-aachen.de/data/}.
\end{itemize}

~

\textbf{Course Management:} Learning Suite

% ------------------------------------------------------------------- software
\subsection*{Software}

We use Python. In class I will use interactive Python notebooks in
\url{https://colab.research.google.com}, which has tradeoffs we will talk about. Your homework can
be completed in any editor or IDE you like. A summary of useful Python commands is at
\url{https://www.pythoncheatsheet.org}.

Course material is on GitHub at
\url{https://github.com/drbob-richardson/stat220/tree/main/F2026}. Slides, code companions,
homework, and midterm practice are posted there unit by unit as we reach them. Please make sure
you have a GitHub account, since you will link it to your name in the first homework.

% --------------------------------------------------------------------- topics
\section*{What We Will Cover}

The course runs in twelve units of roughly three class periods each. The order below is the
starting plan. I may adjust the sequence as we go, and I will announce dates for each unit and
each homework in class and on Learning Suite.

\begin{itemize}[topsep=2pt,itemsep=2pt]
\item \textbf{Statistical Inference.} Signal against noise, standard error, the
  $t$-statistic, $p$-values, the two kinds of error, power, confidence intervals, practical versus
  statistical significance, and the assumptions holding all of it up.
\item \textbf{Regression as Reasoning.} The line, least squares, and what a slope does and does
  not tell you. Inference on a slope, categorical predictors, multiple regression, confounding.
\item \textbf{Building and Trusting a Model.} Many predictors, interactions, collinearity.
  Overfitting, cross-validation, and data leakage. Trees, regularization, variable importance.
\item \textbf{Prediction and Its Uncertainty.} The four layers of prediction error. Prediction
  intervals versus confidence intervals. Extrapolation, regression to the mean, and what happens
  when the world shifts under a model.
\item \textbf{Causal Thinking.} Counterfactuals and randomization. Telling a confounder from a
  collider from a mediator, and why controlling for the wrong one makes things worse.
  Quasi-experiments and difference-in-differences.
\item \textbf{Where the Data Comes From.} Population, sampling frame, and sampling designs.
  Missing data and the mechanisms behind it. Measurement, proxies, and provenance.
\item \textbf{Probability: Randomness, Conditioning, and Bayes.} The rules, independence, and
  what correlated failure does to a system. Conditioning, Bayes' rule, and base rates.
\item \textbf{Random Variables and Distributions.} Random variables, expectation, variance, and
  the common distributions, chosen by reasoning about how the data was generated.
\item \textbf{The Normal Distribution and the Central Limit Theorem.} Densities, the normal,
  $z$-scores. The CLT, the standard error, and how large $n$ actually has to be. Where it fails:
  skew, heavy tails, dependence.
\item \textbf{Estimation, Likelihood, and the Bootstrap.} Estimators, bias and variance,
  shrinkage. Likelihood and what its curvature tells you. The bootstrap and the situations where
  it quietly breaks.
\item \textbf{Categorical Outcomes.} Proportions and rates, absolute versus relative risk.
  Two-way tables, chi-square, Simpson's paradox, logistic regression. Evaluating a classifier:
  thresholds, costs, and calibration.
\item \textbf{Bayesian Reasoning and Decisions.} Updating, and the roles of prior, likelihood,
  and posterior. Choosing priors and testing their influence. Credible intervals, expected loss,
  and the value of information.
\end{itemize}

% ---------------------------------------------------------------------- grades
\section*{Grade Information}

\subsection*{Grade Distribution}
\begin{tabular}{ l l }
Homework & 35\% \\
Midterm Exam & 30\% \\
Final Exam & 35\%
\end{tabular}

\subsection*{Homework}

There is one homework per unit, posted as an incomplete Python notebook. You complete the
questions and upload the finished notebook to Learning Suite.

\textbf{Homework is due at 5:00 PM on the day it is due.} Due dates are announced in class and
posted on Learning Suite. No late work is accepted.

Each homework opens with a simulated data set built so the principle from that unit is visible on
purpose, and where you can check your answer against the truth because you know what generated the
data. It then moves to a real data set and a real question, where nobody knows the truth and the
last part of the problem is usually deciding what can and cannot be said.

Homework is graded on correctness. You may work with other people, but each of you writes your own
answers.

\subsection*{Exams}

Both exams are taken in the \textbf{BYU Testing Center}. Both are closed book and closed computer.
Bring a pencil.

\begin{itemize}[topsep=2pt,itemsep=1pt]
\item \textbf{Midterm: October 19 to October 22.}
\item \textbf{Final: during finals week, December 12 to December 17.}
\end{itemize}

Check \url{https://testing.byu.edu} for hours and for the last time you can start an exam. The
Testing Center is busiest on the final day of any exam window, so do not plan around the last
hour.

The exams are not about software or arithmetic. Every question is a scenario: here is a situation,
here is what somebody concluded, and you say whether the conclusion holds, what would have to be
true for it to hold, or what you would need to know next. Nothing requires a calculation that will
not fit in a margin. This is the same kind of thinking the in-class activities and the last part
of each homework problem are training.

% ------------------------------------------------------------------------- AI
\section*{AI Usage Statement}

I am in the camp of professors who think AI is remarkable and that you should be using it. This
class is a good place to learn where tools like ChatGPT help and where they quietly mislead you.
Debugging with an LLM is far easier than what came before, and I would rather you got good at it
here.

You may use AI on the homework. What you cannot do is skip the understanding, because both exams
are paper and pencil in the Testing Center with no AI and no computer. If you have leaned on a
model to produce answers you could not reconstruct yourself, the exams will show it. To pass this
course you have to demonstrate the reasoning on your own.

\subsection*{When Working In Class}
\begin{enumerate}[topsep=2pt,itemsep=1pt]
\item Look at the course materials and your notes first.
\item If you are stuck, ask a neighbor, a TA, or me.
\item After that, use AI or Google. Use them to understand how you would solve the problem, not
  to be handed the answer.
\end{enumerate}

\subsection*{When Working At Home}
\begin{enumerate}[topsep=2pt,itemsep=1pt]
\item Start with the course notes and materials.
\item If you are still stuck, use Google or an AI tool, or ask on Discord. Both are fine. AI is a
  guide, not a substitute.
\end{enumerate}

\subsection*{Acceptable Uses of AI}
\begin{itemize}[topsep=2pt,itemsep=1pt]
\item Brainstorming approaches to a problem.
\item Clarifying a concept you do not understand.
\item Debugging code and checking syntax.
\item Exploring code style and better practice.
\item Reviewing a statistical method and when it applies.
\end{itemize}

\subsection*{Prohibited Uses of AI}
\begin{itemize}[topsep=2pt,itemsep=1pt]
\item Submitting AI-generated work as your own without editing or understanding it.
\item Using AI to solve a problem for you instead of working through it.
\item Relying on AI so heavily that you cannot reason on paper when the exam comes.
\end{itemize}

\subsection*{Guidelines}
\begin{itemize}[topsep=2pt,itemsep=1pt]
\item Verify and edit AI output for accuracy and clarity.
\item Be able to explain any code or answer you submit.
\item Cite AI contributions if they significantly shaped your work.
\item Remember that the exams are AI-free. If you over-rely on AI now, you will struggle then.
\end{itemize}

\subsection*{Consequences}

Misuse of AI in violation of this policy will be treated as a breach of BYU's academic integrity
standards and may result in grade penalties or further disciplinary action.

% ---------------------------------------------------------------- boilerplate
\section*{Statements on honor code, harassment, and accessibility:}

\textbf{\large Honor Code:}
In keeping with the principles of the BYU Honor Code, students are expected to be honest in all of
their academic work. Academic honesty means, most fundamentally, that any work you present as your
own must in fact be your own work and not that of another. Violations of this principle may result
in a failing grade in the course and additional disciplinary action by the university. Students are
also expected to adhere to the Dress and Grooming Standards. Adherence demonstrates respect for
yourself and others and ensures an effective learning and working environment. It is the
university's expectation, and my own expectation in class, that each student will abide by all
Honor Code standards. Please call the Honor Code Office at 422-2847 if you have questions about
those standards.

\textbf{\large Sexual Harassment:}
Title IX of the Education Amendments of 1972 prohibits sex discrimination against any participant
in an educational program or activity that receives federal funds. The act is intended to eliminate
sex discrimination in education and pertains to admissions, academic and athletic programs, and
university-sponsored activities. Title IX also prohibits sexual harassment of students by
university employees, other students, and visitors to campus. If you encounter sexual harassment or
gender-based discrimination, please talk to your professor or contact one of the following: the
Title IX Coordinator at 801-422-2130; the Honor Code Office at 801-422-2847; the Equal Employment
Office at 801-422-5895; or Ethics Point at \url{http://www.ethicspoint.com}, or 1-888-238-1062
(24-hours).

\textbf{\large Student Disability:}
Brigham Young University is committed to providing a working and learning atmosphere that
reasonably accommodates qualified persons with disabilities. If you have any disability which may
impair your ability to complete this course successfully, please contact the University
Accessibility Center (UAC), 2170 WSC or 422-2767. Reasonable academic accommodations are reviewed
for all students who have qualified, documented disabilities. The UAC can also assess students for
learning, attention, and emotional concerns. Services are coordinated with the student and
instructor by the UAC. If you need assistance or if you feel you have been unlawfully discriminated
against on the basis of disability, you may seek resolution through established grievance policy
and procedures by contacting the Equal Employment Office at 422-5895, D-285 ASB.

\textbf{\large Mental Health Services:}
Barriers to learning are created by stress, anxiety, family and relationship concerns, and personal
crises. If stressful life events or mental health concerns are inhibiting your ability to
participate in daily activities or leading to diminished academic performance, please contact the
BYU Counseling and Psychological Services (CAPS; 1500 WSC, 801-422-3035, \url{caps.byu.edu}). CAPS
provides individual, couples and group counseling to students. These services are confidential and
are provided by the university at no added cost to you. Professional psychologists and counselors
who specialize in helping college students are available 24-hours a day to assist students in
crisis; if you have an emergency during non-business hours (5pm-8am), please contact BYU Police
Dispatch (801-422-2222) who will put you in touch with a counselor.

\end{document}
